\chapter{Cytomegalovirus (CMV)}

\section*{What Is This Test?}
Cytomegalovirus (CMV) testing detects human cytomegalovirus (HHV-5), a ubiquitous double-stranded DNA virus of the \textit{Herpesviridae} family. CMV seroprevalence ranges from 40--50\% in developed countries to >90\% in developing countries. Primary CMV infection in immunocompetent individuals is usually asymptomatic or causes a mild mononucleosis-like syndrome (fever, malaise, lymphadenopathy, mild hepatitis). However, CMV is a major pathogen in two populations: (1) immunocompromised hosts, particularly hematopoietic stem cell transplant (HSCT) and solid organ transplant (SOT) recipients, where CMV reactivation or primary infection causes life-threatening disease including pneumonitis, hepatitis, colitis, esophagitis, retinitis, and encephalitis; and (2) the developing fetus, where congenital CMV infection (transplacental transmission during primary maternal infection) is the leading infectious cause of sensorineural hearing loss, cerebral palsy, and neurodevelopmental disability. CMV serology (IgG, IgM), CMV DNA quantitation by PCR (viral load), CMV antigenemia (pp65 antigenemia assay), and histopathology (tissue biopsy with CMV inclusion bodies) are the principal diagnostic tools. CMV viral load monitoring by quantitative PCR is the standard of care for preemptive therapy and diagnosis of active CMV disease in transplant recipients.

\section*{Alternative Names}
\begin{itemize}
  \item CMV Serology
  \item CMV IgG/IgM
  \item CMV DNA PCR
  \item CMV Viral Load
  \item CMV Antigenemia
  \item CMV pp65
  \item Cytomegalovirus PCR
  \item CMV Culture
  \item CMV Shell Vial Assay
\end{itemize}
\section*{Principle}
CMV serology: CLIA or ELISA detecting CMV-specific IgG and IgM antibodies. CMV IgG indicates past infection (seropositive); CMV IgM suggests recent/active infection but can persist for >12 months after primary infection and may reappear during reactivation. CMV IgG avidity testing helps distinguish primary infection (low avidity) from past infection (high avidity), especially in pregnant women. CMV PCR (viral load): Quantitative real-time PCR detecting CMV DNA in plasma, whole blood, CSF, BAL, or tissue. WHO International Standard (IU/mL) was adopted in 2010 to standardize quantitation across assays, though inter-assay variability persists (up to 0.5 log$_{10}$ difference). Lower limit of quantitation typically 30--200 IU/mL. Results guide preemptive therapy decisions. CMV pp65 antigenemia: Detection of CMV lower matrix phosphoprotein pp65 in peripheral blood leukocytes by immunofluorescence; semi-quantitative (positive cells/200,000 leukocytes). Historically the gold standard, now largely replaced by PCR. CMV shell vial culture: Specimen inoculated onto fibroblast monolayer, centrifuged to enhance infection, and stained after 24--48 hours with fluorescein-labeled anti-CMV immediate-early antigen monoclonal antibody. More rapid than conventional culture (1--2 days vs. 1--4 weeks). Histopathology: Tissue biopsy (GI mucosa, lung) with H\&E stain showing characteristic cytomegalic cells (enlarged cells 25--35 $\mu$m with intranuclear "owl's eye" basophilic inclusion and cytoplasmic inclusions). Confirmed by CMV immunohistochemical staining or \textit{in-situ} hybridization.

\section*{History}
CMV was first described in 1904 by Ribbert in tissue sections from a stillborn infant, who noted the characteristic large inclusion-bearing cells. The virus was first isolated in cell culture by Thomas Weller, Margaret Smith, and Wallace Rowe independently in 1956--1957. The term "cytomegalovirus" was coined in 1960. CMV serologic testing became available in the 1970s, allowing screening of blood donors and transplant donors/recipients. The introduction of ganciclovir (1988), valganciclovir (2001), foscarnet, and cidofovir provided effective therapy but with significant toxicity. The pp65 antigenemia assay was developed in the 1990s as a semi-quantitative marker of active CMV infection in transplant recipients. CMV DNA PCR became widely adopted in the 2000s, and the WHO International Standard for CMV DNA quantitation was established in 2010 to harmonize results across laboratories. The CMV T-cell immunity assays (QuantiFERON-CMV, CMV ELISPOT) were developed in the 2010s to assess CMV-specific cell-mediated immunity and predict risk of CMV reactivation post-transplant. Letermovir, a novel CMV terminase inhibitor, was FDA-approved in 2017 for CMV prophylaxis in HSCT recipients.

\section*{Why This Test Is Ordered}
\begin{itemize}
  \item Pre-transplant CMV serologic screening of donors and recipients (D/R serostatus): The most important predictor of post-transplant CMV risk. High-risk: D+/R-; intermediate: D+/R+, D-/R+ (SOT). Guiding prophylaxis versus preemptive therapy strategy \cite{kotton2018third}
  \item Weekly CMV viral load monitoring in high-risk transplant recipients (HSCT, SOT) for the first 100 days post-transplant (or longer) as part of preemptive therapy strategy: rising CMV DNA or exceeding a center-specific threshold triggers initiation of antiviral therapy to prevent CMV end-organ disease \cite{boeckh2009how}
  \item Diagnosis of CMV end-organ disease (pneumonitis, colitis, hepatitis, retinitis, encephalitis) in immunocompromised patients with compatible clinical and radiographic findings plus detection of CMV by PCR, antigenemia, and/or tissue biopsy
  \item Evaluation of primary CMV infection in immunocompetent adults with mononucleosis-like syndrome (fever, lymphadenopathy, atypical lymphocytosis) when EBV testing (Monospot, EBV serology) is negative
  \item Diagnosis of congenital CMV infection: PCR of neonatal urine (within the first 2--3 weeks of life, as detection after 3 weeks may represent postnatal acquisition), saliva, or blood \cite{rawlinson2017congenital}. Dried blood spot (DBS) PCR from newborn screening cards is used for retrospective diagnosis. CMV serology in the pregnant woman with mononucleosis-like illness, fetal ultrasound abnormalities (intracranial calcifications, microcephaly, hepatosplenomegaly, echogenic bowel), or known CMV exposure; IgM, IgG, and IgG avidity testing are essential to estimate timing of maternal infection \cite{lazzarotto2008new}
  \item Screening of blood and organ donors for CMV IgG to prevent transfusion- and donor-derived CMV transmission to CMV-seronegative recipients
  \item CMV retinitis screening in HIV/AIDS patients with CD4 <50 cells/$\mu$L: indirect ophthalmoscopy; CMV PCR of vitreous fluid may be confirmatory
\end{itemize}

\section*{Primary vs Secondary Test}
CMV serology (IgG) is the primary test for pre-transplant donor/recipient screening. CMV quantitative PCR (viral load) is the primary test for monitoring and diagnosis of active CMV infection in transplant recipients. Tissue biopsy with histopathology remains the gold standard for the diagnosis of CMV end-organ disease, particularly for GI disease (colitis, esophagitis). CMV serology (IgM, IgG avidity) is the primary test for evaluation of primary maternal infection in pregnancy.

\section*{How This Test Is Performed}
Serology: Venous blood in SST tube; serum separated and refrigerated or frozen. CMV PCR: EDTA whole blood or plasma; plasma is preferred for viral load quantitation. Specimens should be centrifuged and plasma separated within 6 hours. Quantitative results in IU/mL. Urine CMV PCR: Clean-catch neonatal urine; the highest viral loads are in urine. CSF CMV PCR: 0.5--1 mL CSF in sterile tube. Tissue biopsy: Formalin-fixed paraffin-embedded tissue for histopathology with CMV immunohistochemistry. Bronchoalveolar lavage (BAL) fluid for CMV PCR and shell vial culture. Vitreous fluid: Collected by ophthalmologist for CMV PCR when retinitis is suspected.

\section*{What Preparation Is Needed}
No special preparation for blood, urine, or BAL collection. For tissue biopsy, standard preoperative preparation applies.

\section*{Contraindications and Precautions}
\begin{itemize}
  \item CMV IgM can persist for >12 months after primary infection and may be detected during CMV reactivation or in response to other herpesvirus infections (EBV, HHV-6) due to cross-reactivity or polyclonal B-cell stimulation. CMV IgM alone is insufficient to diagnose primary CMV infection; IgG avidity testing or paired serology demonstrating seroconversion (IgG negative to positive) is more specific. In transplant recipients with CMV seropositivity, CMV IgG testing post-transplant is not useful for diagnosing acute disease; viral load (PCR) and/or antigenemia are the appropriate tests. Quantitative CMV PCR values can vary significantly between assays and laboratories (up to 0.5 log$_{10}$), so serial monitoring should ideally be performed using the same assay and laboratory. Plasma PCR may be less sensitive than whole blood PCR for detecting CMV reactivation (CMV DNA is predominantly cell-associated in leukocytes), but plasma PCR is better standardized. The definition of an actionable CMV viral load for preemptive therapy varies by center and patient population; no universal threshold exists. Rising viral load or exceeding a local threshold (typically 500--5,000 IU/mL in plasma) triggers preemptive therapy.
\end{itemize}
\section*{Risks}
Standard venipuncture, lumbar puncture, and biopsy risks apply. No risks from the serologic or molecular tests themselves. The major clinical risk is CMV disease in an immunocompromised patient who is not adequately monitored or treated: CMV pneumonitis in HSCT recipients carries 50--90\% mortality without treatment, reduced to 10--30\% with ganciclovir plus IVIG \cite{ljungman2019guidelines}.

\section*{What Happens After the Test}
CMV serology: results in 1--24 hours. CMV PCR: results in 24--48 hours. CMV shell vial culture: 24--48 hours. Histopathology: 2--5 days. Rising CMV viral load in a transplant recipient triggers preemptive antiviral therapy (valganciclovir 900 mg PO BID, or IV ganciclovir 5 mg/kg q12h) before end-organ disease develops. Established CMV disease requires treatment-dose antivirals for 14--21 days with close monitoring of CBC (ganciclovir causes myelosuppression/neutropenia). In congenital CMV infection, oral valganciclovir for 6 months (initiated within the first month of life) has been shown to improve hearing and neurodevelopmental outcomes in symptomatic infants.

\section*{Understanding Results}

\subsection*{Below Normal}
\begin{itemize}
  \item \textbf{CMV IgG negative (seronegative):} No prior CMV infection. In transplant: D-/R- status for both donor and recipient = low risk for CMV disease, provided CMV-seronegative or leukoreduced blood products are used. In pregnancy: Risk of primary CMV infection exists; counsel about CMV prevention (hand hygiene, avoid sharing food/utensils with young children who are common sources of CMV transmission)
  \item \textbf{CMV PCR not detected or <LLOQ:} No detectable CMV DNA. In transplant recipients on preemptive therapy monitoring, indicates no CMV reactivation. In symptomatic patients, makes active CMV disease unlikely (though GI CMV disease may have negative plasma PCR and be positive only on tissue PCR or histopathology)
  \item \textbf{CMV culture negative:} No replicating virus detected. PCR is more sensitive than culture
\end{itemize}

\subsection*{Above Normal}
\begin{itemize}
  \item \textbf{CMV IgG positive (seropositive):} Past CMV infection. Does NOT distinguish between latent and active infection; most CMV-seropositive transplant recipients are at risk for reactivation when immunosuppressed. In pregnancy: CMV IgG positive with negative IgM and high IgG avidity indicates remote (past) infection, not primary infection during pregnancy, with a very low risk (<1\%) of fetal transmission. If IgM is positive and IgG avidity is low, primary infection during pregnancy is likely (30--40\% risk of fetal transmission)
  \item \textbf{CMV IgM positive:} Suggestive of recent/active CMV infection. Not specific for primary infection (can be positive during reactivation, co-infection with EBV, or due to cross-reactivity). Requires further evaluation with IgG avidity (or paired serology) to assess timing of infection
  \item \textbf{CMV IgG low avidity:} Indicates primary infection within the past 3--4 months. High clinical relevance in pregnancy: low avidity IgG combined with IgM positivity strongly suggests primary CMV infection during pregnancy with high risk (~30--40\%) of fetal transmission
  \item \textbf{CMV DNA detected (viral load quantitated in IU/mL):} Indicates active CMV replication. In transplant recipients: rising viral load or load above the institution's threshold triggers preemptive antiviral therapy. In a symptomatic immunocompromised patient, a high viral load (>10,000 IU/mL) supports the diagnosis of CMV disease. Viral load trends (rising, declining) guide duration of therapy and monitoring
  \item \textbf{CMV shell vial culture positive (fluorescent focus detected):} Confirms replicating CMV. The rapid culture method significantly shortened the time to detection from weeks to 24--48 hours
  \item \textbf{Tissue histopathology positive (cytomegalic cells with owl's eye inclusions and CMV IHC positive):} Gold standard for diagnosing CMV end-organ disease (CMV colitis, esophagitis, pneumonitis). Histopathology plus detection of CMV DNA or culture of the same specimen confirms invasive disease
\end{itemize}

\section*{Normal Results}
CMV IgG: \textbf{Negative (seronegative)} or \textbf{Positive (seropositive, latent)}. CMV IgM: \textbf{Negative}. CMV PCR: \textbf{Not detected} (<LLOQ). CMV urine culture/PCR (neonate): \textbf{Negative} (within first 2--3 weeks of life).

\section*{Follow-up Tests}
\begin{itemize}
  \item \textbf{CMV viral load (quantitative PCR) serial monitoring:} Weekly for the first 100 days post-transplant (HSCT) or per institutional protocol for SOT. A rising trend or exceeding the preemptive therapy threshold triggers treatment
  \item \textbf{CMV IgG avidity testing:} When CMV IgM is positive in a pregnant woman, avidity determines the likelihood of recent primary infection and the risk of congenital CMV. Low avidity = primary infection within the past 3--4 months; high avidity = remote infection
  \item \textbf{Amniocentesis with CMV PCR of amniotic fluid:} For pregnant women with confirmed primary CMV infection; performed at $\geq$21 weeks gestation, $\geq$6 weeks after maternal infection; positive PCR confirms fetal infection
  \item \textbf{Neonatal urine/saliva CMV PCR:} The diagnostic test for congenital CMV infection; must be performed within the first 2--3 weeks of life to differentiate congenital from postnatal acquisition
  \item \textbf{CMV antiviral resistance testing (genotypic: UL97 kinase and UL54 DNA polymerase sequencing):} When CMV viral load rises or fails to decline after >2 weeks of appropriate-dose ganciclovir therapy, suggesting drug resistance. UL97 mutations confer ganciclovir resistance; UL54 mutations confer cross-resistance to cidofovir and/or foscarnet
  \item \textbf{Tissue biopsy with CMV IHC:} For suspected CMV end-organ disease, particularly GI disease (colitis, esophagitis) where plasma PCR may be negative despite invasive GI disease
\end{itemize}

\section*{References}
\bibliographystyle{unsrtnat}
\bibliography{references}

\clearpage
\section*{Regulatory Status and Authorization Date}
This chapter describes a test category, not a specific commercial product. FDA, CDSCO, EU IVDR, and MHRA approval or registration dates therefore cannot be assigned without the assay manufacturer, product name, instrument, and intended use. For laboratory-developed tests, an individual agency approval date may not exist. See the Preface for the interpretation of regulatory status.

\clearpage
